%% bare_jrnl.tex
%% V1.4b
%% 2015/08/26
%% by Michael Shell
%% see http://www.michaelshell.org/
%% for current contact information.
%%
%% This is a skeleton file demonstrating the use of IEEEtran.cls
%% (requires IEEEtran.cls version 1.8b or later) with an IEEE
%% journal paper.
%%
%% Support sites:
%% http://www.michaelshell.org/tex/ieeetran/
%% http://www.ctan.org/pkg/ieeetran
%% and
%% http://www.ieee.org/

%%*************************************************************************
%% Legal Notice:
%% This code is offered as-is without any warranty either expressed or
%% implied; without even the implied warranty of MERCHANTABILITY or
%% FITNESS FOR A PARTICULAR PURPOSE! 
%% User assumes all risk.
%% In no event shall the IEEE or any contributor to this code be liable for
%% any damages or losses, including, but not limited to, incidental,
%% consequential, or any other damages, resulting from the use or misuse
%% of any information contained here.
%%
%% All comments are the opinions of their respective authors and are not
%% necessarily endorsed by the IEEE.
%%
%% This work is distributed under the LaTeX Project Public License (LPPL)
%% ( http://www.latex-project.org/ ) version 1.3, and may be freely used,
%% distributed and modified. A copy of the LPPL, version 1.3, is included
%% in the base LaTeX documentation of all distributions of LaTeX released
%% 2003/12/01 or later.
%% Retain all contribution notices and credits.
%% ** Modified files should be clearly indicated as such, including  **
%% ** renaming them and changing author support contact information. **
%%*************************************************************************


% *** Authors should verify (and, if needed, correct) their LaTeX system  ***
% *** with the testflow diagnostic prior to trusting their LaTeX platform ***
% *** with production work. The IEEE's font choices and paper sizes can   ***
% *** trigger bugs that do not appear when using other class files.       ***                          ***
% The testflow support page is at:
% http://www.michaelshell.org/tex/testflow/



\documentclass[journal]{IEEEtran}
%
% If IEEEtran.cls has not been installed into the LaTeX system files,
% manually specify the path to it like:
% \documentclass[journal]{../sty/IEEEtran}





% Some very useful LaTeX packages include:
% (uncomment the ones you want to load)


% *** MISC UTILITY PACKAGES ***
%
%\usepackage{ifpdf}
% Heiko Oberdiek's ifpdf.sty is very useful if you need conditional
% compilation based on whether the output is pdf or dvi.
% usage:
% \ifpdf
%   % pdf code
% \else
%   % dvi code
% \fi
% The latest version of ifpdf.sty can be obtained from:
% http://www.ctan.org/pkg/ifpdf
% Also, note that IEEEtran.cls V1.7 and later provides a builtin
% \ifCLASSINFOpdf conditional that works the same way.
% When switching from latex to pdflatex and vice-versa, the compiler may
% have to be run twice to clear warning/error messages.






% *** CITATION PACKAGES ***
%
%\usepackage{cite}
% cite.sty was written by Donald Arseneau
% V1.6 and later of IEEEtran pre-defines the format of the cite.sty package
% \cite{} output to follow that of the IEEE. Loading the cite package will
% result in citation numbers being automatically sorted and properly
% "compressed/ranged". e.g., [1], [9], [2], [7], [5], [6] without using
% cite.sty will become [1], [2], [5]--[7], [9] using cite.sty. cite.sty's
% \cite will automatically add leading space, if needed. Use cite.sty's
% noadjust option (cite.sty V3.8 and later) if you want to turn this off
% such as if a citation ever needs to be enclosed in parenthesis.
% cite.sty is already installed on most LaTeX systems. Be sure and use
% version 5.0 (2009-03-20) and later if using hyperref.sty.
% The latest version can be obtained at:
% http://www.ctan.org/pkg/cite
% The documentation is contained in the cite.sty file itself.

% *** graphics and diagram packages ***

\usepackage{graphicx}
\usepackage{tikz}

\usetikzlibrary{
    arrows.meta,
    positioning,
    shapes.geometric
}

\newcommand{\scenenumber}[1]{%
    {\scriptsize\textit{#1}}%
}


% correct bad hyphenation here
\hyphenation{op-tical net-works semi-conduc-tor}


\begin{document}
%
% paper title
% Titles are generally capitalized except for words such as a, an, and, as,
% at, but, by, for, in, nor, of, on, or, the, to and up, which are usually
% not capitalized unless they are the first or last word of the title.
% Linebreaks \\ can be used within to get better formatting as desired.
% Do not put math or special symbols in the title.
\title{Fleet Progress Report}
\author{by Taylor Belcher \\ CSCE552 Game Development, Dr. JJ Shepherd}
\maketitle

\tableofcontents 

\section{General Progress Update}

Summer camps at GSSM have now ended and I am done with the AI in Education conference that I was a speaker at. \emph{I have no other obligations the next two weeks except making the game for this course.} I anticipate that I will still accomplish all necessary tasks by their deadlines. A list of tasks I did this week is below. Project commit history at https://github.com/hotdogontology/CSCE552/commits/main/
\begin{enumerate}
    \item Completed tutorial on GDScript.
    \item Completed tutorial on making 2D games in Godot.
    \item Made stand-in ship model in Blockbench.
    \item Started and abandoned tutorial on making 3D games in Godot in favor of beginning work on game.
\end{enumerate}



\section{Features Currently Implemented}

I have been able to create a basic 3D scene, import my model, and have it move around with the camera in approximately the correct place. It is very much NOT a game right now, but I am happy with my progress given that I have never made a 3D game before or used Godot and my other obligations so far.

\begin{figure}[!t]
    \centering
    \includegraphics[width=\columnwidth]{gameplay.png}
    \caption{screenshot of current game}
    \label{fig:game-play}
\end{figure}

\begin{figure}[!t]
    \centering
    \includegraphics[width=\columnwidth]{maineditor.png}
    \caption{current 3d editor view of main scene}
    \label{fig:editor-godot}
\end{figure}


\begin{enumerate}
    \item A basic scene tree has been established in the project. Current scenes are
    \begin{enumerate}
        \item Main
        \item Player
        \item Trench
        \item TrenchScraper (Building)
        \end{enumerate}
    \item A player ship is drawn within a trench scene.
    \item The player can move the ship in four directions.
\end{enumerate}
\section{Features To Be Implemented}
All other features as described in the GDD submitted last week. Essentially, now that I understand Godot a little, I need to actually make the game.
\begin{enumerate}
    \item Smoother player movement that feels like a ship flying
    \item Shooting weapons
    \item Collisions with objects and enemy weapons
    \item Health and shield systems
    \item Loadout scene and systems
    \item Enemy and obstacle spawn and movement
    \item Random level generation and on-rail levels
    \item Galaxy map scene
    \item All menus
\end{enumerate}
\section{Features Removed}

I still plan to make the game as described in the GDD to at least the minimum specifications.
\section{Problems Encountered}

\begin{enumerate}
    \item The largest obstacle for me this past week was time. I had obligations for the AI conference from 9 AM - 5 PM Monday - Thursday and had to work in the evenings. \textbf{But now I have the next two weeks completely free.}
    \item Audacity is audio-editing software, but as far as I can tell, it cannot allow me to write music. I did not realize this. I will need to download additional tracking software. I am looking into Reaper and NES extensions to write some chiptune music as well as investigating if I can export sounds or music from PICO-8's built-in tracker. 
    \item My triangle spaceship model is missing coloring on some faces and I am not sure why. It was my first blockbench model ever, so I probably just colored the mesh incorrectly. The triangle model is a placeholder for the final model, so this is mainly a non-issue. 
    \item Related to the above, making 3D art in Blockbench is very difficult! But I am sticking to basic shapes to manage.
    \item Godot makes scene management and built-in features as described in lecture very manageable, but I am finding that style of organization for writing a game far more un-intuitive than coding everything like I have in PICO-8 games. My sense of ``\emph{here is what code you should write to implement that feature}" that I have developed is often unrelated to the new sense of ``\emph{here is the node type that you should create to implement that feature and what characteristics you should edit}" that Godot requires. That is, the learning curve is steeper than I anticipated. However, the tutorials helped and I am getting the hang of it. 
    \item I can get the player to move up, down, left, right, but the movement does not \emph{feel} like a spaceship moving yet. I need to implement the nose tipping up and down and the ship banking. I also need to have the camera probably move slightly with the ship. I will need to do some research into how to get a better game feel for player controls.
\end{enumerate}

\section{Team Tasks}

\begin{itemize}

    \item \textbf{Taylor, July 17--20:}
    Improve player movement, add collisions, enemies, obstacles, weapons, and loadout system.

    \item \textbf{Taylor, July 20--22:}
    Create MWE in trench tutorial scene with randomly generated enemies. Create at least one on-rails hand-crafted level to test the PathFollow node. 

    \item \textbf{Taylor, July 22--24:}
    Create set-up for all randomly generated scenes, create galaxy map and mission selection system.

    \item \textbf{Taylor, July 24--29:}
    Create as many on-rails levels as possible. 

    \item \textbf{Taylor, July 29--31:}
    Finishing details, playtest and revise, submit game.

    \item \textbf{Clayton, July 20--31:}
    Create art and music as necessary to assist Taylor. Playtest and give feedback.

\end{itemize}




\section{Programming Tasks}

\begin{itemize}

    \item \textbf{Taylor, July 17--20:}
    Improved player movement, collisions signals, enemy movement and behavior, build loadout system.

    \item \textbf{Taylor, July 20--22:}
    Generating enemies and obstacles for randomly generated levels. Figure out how the path follow node interacts with player input.  

    \item \textbf{Taylor, July 22--24:}
    Mission select and galaxy map logic.

    \item \textbf{Taylor, July 24--29:}
    Create as many on-rails levels as possible. 

    \item \textbf{Taylor, July 29--31:}
    Finishing details, playtest and revise, submit game.

\end{itemize}

\section{Art Tasks}

Minimum art assets we need to create. These are all assigned to Taylor by default. Clayton will assist where he can.

\begin{itemize}
    \item \textbf{Taylor, July 17--22:}
    \begin{itemize}
    \item Enemy Spaceship 1 Model
    \item Enemy Spaceship 2 Model
    \item Enemy Spaceship 3 Model
    \item Enemy Turret 1 Model
    \item Laser
    \item Bomb
    \item Building 1 Model
    \item Asteroid Model
    \item Asteroid Field Icon
    \end{itemize}
    \item \textbf{Taylor, July 22--24:}
    \begin{itemize}
    \item Space Station Model
    \item Space Station Icon
    \item Planet Icon
    \item Planet Background
    \item Improved Player Spaceship Model
    \item Carrier Ship Model
    \item Battery Loadout Icon
    \item Shield Generator Loadout Icon
    \item Missile Loadout Icon
    \item Loadout Bay Puzzle
    \item Galaxy Map
    \end{itemize}

\end{itemize}

\section{Testing Tasks}

\begin{itemize}

    \item \textbf{Taylor, July 17--20:}
    Test movement, firing, collisions, energy, shields, armor, pickups, and loadout validation as they are implemented.

    \item \textbf{Taylor, July 20--22:}
    Test all scene transitions and verify that fleet and campaign data persist.

    \item \textbf{Taylor, July 22--24:}
    Test victory, failure, retreat, replacement-fighter, save, load, and game-end conditions.

    \item \textbf{Taylor, July 24--29:}
    Conduct playtests focused on controls, objective clarity, interface readability, mission length, and difficulty. 
    
    \item \textbf{Taylor, July 29--31:}
    Perform a final test and correct crashes, soft locks, and
    progression blockers.

    \item \textbf{Clayton, July 20--31:}
    Conduct playtests focused on controls, objective clarity, interface readability, mission length, and difficulty.
    

\end{itemize}

\end{document}